\chapter{章节文本}

\section{文本}

\subsection{短篇文本}

庆历四年春，滕子京谪守巴陵郡。越明年，政通人和，百废具兴，乃重修岳阳楼，增其旧制，刻唐贤今人诗赋于其上，属予作文以记之。

予观夫巴陵胜状，在洞庭一湖。衔远山，吞长江，浩浩汤汤，横无际涯，朝晖夕阴，气象万千，此则岳阳楼之大观也，前人之述备矣。然则北通巫峡，南极潇湘，迁客骚人，多会于此，览物之情，得无异乎？

若夫淫雨霏霏，连月不开，阴风怒号，浊浪排空，日星隐曜，山岳潜形，商旅不行，樯倾楫摧，薄暮冥冥，虎啸猿啼。登斯楼也，则有去国怀乡，忧谗畏讥，满目萧然，感极而悲者矣。

至若春和景明，波澜不惊，上下天光，一碧万顷，沙鸥翔集，锦鳞游泳，岸芷汀兰，郁郁青青。而或长烟一空，皓月千里，浮光跃金，静影沉璧，渔歌互答，此乐何极！登斯楼也，则有心旷神怡，宠辱偕忘，把酒临风，其喜洋洋者矣。

嗟夫！予尝求古仁人之心，或异二者之为，何哉？不以物喜，不以己悲，居庙堂之高则忧其民，处江湖之远则忧其君。是进亦忧，退亦忧。然则何时而乐耶？其必曰：“先天下之忧而忧，后天下之乐而乐”乎！噫！微斯人，吾谁与归？

时六年九月十五日。

\subsection{长篇文本}

奇点炸弹脱离导轨后，沿一条由母舰发出的力场束加速，直奔目标恒星。过了不长的一段时间，这颗灰尘似的黑洞高速射入了恒星表面火的海洋。想象在太平洋的中部突然出现一个半径一百公里的深井，就可以大概把握这时的情形。巨量的恒星物质开始被吸入黑洞，那汹涌的物质洪流从所有方向会聚到一点并消失在那里，物质吸入时产生的辐射在恒星表面产生一团剌目的光球，仿佛恒星戴上了一个光彩夺目的钻石戒指。随着黑洞向恒星内部沉下去，光团暗淡下来，可以看到它处于一个直径达几百万公里的大旋涡正中，那巨大的旋涡散射着光团的强光，缓缓转动着，呈现出飞速变幻的色彩，使恒星从这个方向看去仿佛是一张狰狞的巨脸。很快，光团消失了，旋涡渐渐消失，恒星表面似乎又恢复了它原来的色彩和光度。但这只是毁灭前最后的平静，随着黑洞向恒星中心下沉，这个贪婪的饕餐者更疯狂地吞食周围密度急剧增高的物质，它在一秒钟内吸入的恒星物质总量可能有上百个中等行星。黑洞巨量吸入时产生的超强辐射向恒星表面漫延，由于恒星物质的阻滞，只有一小部分到达了表面，但其余的辐射把它们的能量留在了恒星内部，这能量快速破坏着恒星的每一个细胞，从整体上把它飞快地拉离平衡态。从外部看，恒星的色彩在缓缓变化，由浅红色变为明黄色，从明黄色变为鲜艳的绿色，从绿色变为如洗的碧蓝，从碧蓝变为恐怖的紫色。这时，在恒星中心的黑洞产生的辐射能已远远大于恒星本身辐射的能量，随着更多的能量以非可见光形式溢出恒星，这紫色在加深加深，这颗恒星看上去象太空中一个在忍受着超级痛苦的灵魂，这痛苦在急剧增大，紫色已深到了极限，这颗恒星用不到一个小时的时间走完了它未来几十亿年的旅程。

一团似乎吞没整个宇宙的强光闪起，然后慢慢消失，在原来恒星所在的位置上，可以看到一个急剧膨涨的薄球层，象一个被吹大的气球，这是被炸飞的恒星表面。随着薄球层体积的增大，它变得透明了，可以看到它内部的第二个膨涨的薄球层，然后又可以看到更深处的第三个薄球层……这个爆炸中的恒星，就象宇宙中突然显现的一个套一个的一组玲笼剔透的缕花玻璃球，其中最深处的一个薄球层的体积也是恒星原来体积的几十万倍。

当爆炸的恒星的第一层膨涨外壳穿过那个橙黄色行星时，它立刻被汽化了。其实在这整个爆炸的壮丽场景中根本就看不到它，同那膨涨的恒星外壳相比，它只是一粒微不足道的灰尘，其大小甚至不能成为那几层缕花玻璃球上的一个小点。

\subsection{段落}

\subsubsection{测试段落（paragraph）的效果}

\paragraph{立法目的} 
为了保护民事主体的合法权益，调整民事关系，维护社会和经济秩序，适应中国特色社会主义发展要求，弘扬社会主义核心价值观，根据宪法，制定本法。

\paragraph{调整范围} 
民法调整平等主体的自然人、法人和非法人组织之间的人身关系和财产关系。

\subsubsection{测试子段落（subparagraph）的效果}

\subparagraph{平等原则} 
民事主体在民事活动中的法律地位一律平等。

\subparagraph{自愿原则} 
民事主体从事民事活动，应当遵循自愿原则，按照自己的意思设立、变更、终止民事法律关系。

\subparagraph{公平原则}
民事主体从事民事活动，应当遵循公平原则，合理确定各方的权利和义务。

\subparagraph{诚信原则}
民事主体从事民事活动，应当遵循诚信原则，秉持诚实，恪守承诺。

\section{列表}

\subsection{有序列表}
\begin{itemize}
    \item 水星（Mercury）
    \item 金星（Venus）
    \item 地球（Earth）
    \begin{itemize}
        \item 我们在这里。
    \end{itemize}
    \item 火星（Mars）
    \item 木星（Jupiter）
    \item 土星（Saturn）
    \item 天王星（Uranus）
    \item 海王星（Neptune）
\end{itemize}

\subsection{无序列表}
\begin{enumerate}
    \item 氢（H）
    \item 氦（He）
    \item 锂（Li）
    \item 铍（Be）
    \item 硼（B）
    \item 碳（C）
    \item 氮（N）
    \item 氧（O）
    \begin{enumerate}
        \item 生命之源。
    \end{enumerate}
    \item 氟（F）
    \item 氖（Ne）
\end{enumerate}

\section{字体}

\subsection{中文字体}

\subsubsection{常规}
衬线：{\rmfamily 一只敏捷的棕色狐狸跳过一只懒惰的狗。}

等线：{\sffamily 一只敏捷的棕色狐狸跳过一只懒惰的狗。}

等宽：{\ttfamily 一只敏捷的棕色狐狸跳过一只懒惰的狗。}

\subsubsection{粗体}
衬线：{\rmfamily\bfseries 一只敏捷的棕色狐狸跳过一只懒惰的狗。}

等线：{\sffamily\bfseries 一只敏捷的棕色狐狸跳过一只懒惰的狗。}

等宽：{\ttfamily\bfseries 一只敏捷的棕色狐狸跳过一只懒惰的狗。}

\subsubsection{斜体}
衬线：{\rmfamily\itshape 一只敏捷的棕色狐狸跳过一只懒惰的狗。}

等线：{\sffamily\itshape 一只敏捷的棕色狐狸跳过一只懒惰的狗。}

等宽：{\ttfamily\itshape 一只敏捷的棕色狐狸跳过一只懒惰的狗。}

\subsection{英文字体}

\subsubsection{常规}
衬线：{\rmfamily The quick brown fox jumps a lazy dog.}

等线：{\sffamily The quick brown fox jumps a lazy dog.}

等宽：{\ttfamily The quick brown fox jumps a lazy dog.}

\subsubsection{粗体}
衬线：{\rmfamily\bfseries The quick brown fox jumps a lazy dog.}

等线：{\sffamily\bfseries The quick brown fox jumps a lazy dog.}

等宽：{\ttfamily\bfseries The quick brown fox jumps a lazy dog.}

\subsubsection{斜体}
衬线：{\rmfamily\itshape The quick brown fox jumps a lazy dog.}

等线：{\sffamily\itshape The quick brown fox jumps a lazy dog.}

等宽：{\ttfamily\itshape The quick brown fox jumps a lazy dog.}

\subsection{字号}
{\Huge Size Huge 字号}

{\huge Size huge 字号}

{\LARGE Size LARGE 字号}

{\Large Size Large 字号}

{\large Size large 字号}

{\normalsize Size normalsize 字号}

{\small Size small 字号}

{\footnotesize Size footnotesize 字号}

{\scriptsize Size scriptsize 字号}

{\tiny Size tiny 字号}

\section{引用}

\subsection{单一引用}
\cref{part:笔记模板测试}

\cref{chap:章节文本}

\cref{sec:文本}

\cref{subsec:中文字体}

\cref{subsubsec:测试子段落（subparagraph）的效果}

\subsection{多个引用}
\cref{sec:文本}

\cref{sec:文本,sec:列表,sec:字体}

\cref{sec:文本,sec:列表,sec:gamma}

\cref{sec:文本,sec:列表}

\section[Γ 函数]{$\Gamma$函数}[gamma]

测试当一个章节标题包含数学公式，需要和标题名不同的标签名和PDF大纲中的标题名。

这是\Cref{sec:gamma}吗？